\documentclass[11pt]{article}
\usepackage[a4paper,margin=1in]{geometry}
\usepackage{fontspec}
\usepackage[boldfont=false]{xeCJK}
\setCJKmainfont{FandolSong-Regular.otf}
\usepackage{amsmath}
\usepackage{amssymb}
\usepackage{array}
\usepackage{booktabs}
\usepackage{graphicx}
\usepackage{adjustbox}
\usepackage{enumitem}
\setlist[itemize]{topsep=2pt,partopsep=0pt,parsep=0pt,itemsep=2pt,leftmargin=1.4em}
\setlist[enumerate]{topsep=2pt,partopsep=0pt,parsep=0pt,itemsep=2pt,leftmargin=1.6em}
\usepackage{parskip}
\usepackage{fvextra}
\fvset{breaklines=true,breakanywhere=true,fontsize=\small}
\setlength{\parindent}{0pt}

\begin{document}

\begin{center}
  {\LARGE\bfseries Temperature}\\[0.35em]
  {\large A Level Physics}
\end{center}
\vspace{0.6em}

\section*{Thermal equilibrium}

\textbf{Heat} 热量 (thermal energy) flows from a \textbf{higher temperature} to a \textbf{lower temperature}. When two bodies touch, energy moves until their temperatures are equal — they reach \textbf{thermal equilibrium} 热平衡. At equilibrium there is no net flow of energy.

Two regions at the \textbf{same temperature} 温度 are in thermal equilibrium with each other — no net energy flows, even though particles still exchange energy.

Temperature decides the \textbf{direction} of heat flow. It is not a measure of how much \textbf{thermal energy} 热能 a body holds. A small cup of boiling water (100 °C) holds far less energy than a swimming pool at 25 °C, but a piece of metal put in the cup gains energy while one put in the pool loses it.

\section*{Measuring temperature}

Any physical property that \textbf{changes in a repeatable way with temperature} can make a \textbf{thermometer} 温度计. Examples:

\begin{itemize}
\item \textbf{density of a liquid} — a liquid-in-glass thermometer (mercury or alcohol). As the temperature rises, the liquid expands and rises up a narrow \textbf{capillary} 毛细管.

\item \textbf{volume of a gas at constant pressure} — a gas thermometer. The gas volume rises in step with the absolute temperature.

\item \textbf{resistance of a metal} — a resistance thermometer. A metal's \textbf{resistance} 电阻 rises nearly in step with temperature over a wide range.

\item \textbf{e.m.f. of a thermocouple} — a \textbf{thermocouple} 热电偶 is two different metals joined at two points; the \textbf{electromotive force} 电动势 it makes depends on the temperature difference between the joins.

\end{itemize}

Different thermometers can read slightly differently if the property does not change in a straight line; they agree only at the \textbf{calibration} 校准 points.

\par\medskip
\noindent\begin{minipage}{\linewidth}
\centering
\adjustbox{max width=\linewidth, max totalheight=0.42\textheight}{%
  \includegraphics{/Users/joschka/Desktop/Programming/eXercise/output/eXam/handouts/a_level_physics/assets/14-p261-gas-thermometer.pdf}}\\[0.3em]
\emph{A constant-volume gas thermometer — the gas pressure is found from the height difference h}
\end{minipage}\par\medskip

\section*{Thermodynamic temperature scale}

The \textbf{thermodynamic temperature} 热力学温度 (or \textbf{absolute temperature} 绝对温度) scale does not depend on any one substance — only on the laws of thermodynamics. Its unit is the \textbf{kelvin} 开尔文 (K).

\subsection*{Absolute zero}

The lowest possible temperature is \textbf{zero kelvin} ($0\ \text{K}$), called \textbf{absolute zero} 绝对零度. There a system has its least possible \textbf{internal energy} 内能 — particles have no random motion to speak of. Nothing can be cooled below this.

\par\medskip
\noindent\begin{minipage}{\linewidth}
\centering
\adjustbox{max width=\linewidth, max totalheight=0.42\textheight}{%
  \includegraphics{/Users/joschka/Desktop/Programming/eXercise/output/eXam/handouts/a_level_physics/assets/14-p262-absolute-zero-extrapolation.pdf}}\\[0.3em]
\emph{Extrapolating the pressure–temperature line back to zero pressure gives absolute zero, about −273 °C}
\end{minipage}\par\medskip

\subsection*{Celsius scale}

The \textbf{Celsius} 摄氏度 scale $\theta$ is shifted from the thermodynamic scale by a fixed amount:

\[
T / \text{K} = \theta / {}^{\circ}\text{C} + 273.15.
\]

So $0\ ^{\circ}\text{C} = 273.15\ \text{K}$ and $100\ ^{\circ}\text{C} = 373.15\ \text{K}$. A kelvin and a degree Celsius are the \textbf{same size}, so a temperature \textbf{difference} of $1\ \text{K}$ equals $1\ ^{\circ}\text{C}$ — but the absolute values differ by $273.15$.

In gas-law calculations you must always use \textbf{absolute} temperatures in kelvin. Using °C gives wrong answers.

\section*{Specific heat capacity}

The \textbf{specific heat capacity} 比热容 $c$ of a substance is the \textbf{energy} 能量 needed to raise the temperature of \textbf{unit mass} by \textbf{one kelvin}:

\[
c = \frac{Q}{m \Delta T} \qquad\Longleftrightarrow\qquad Q = m c \Delta T.
\]

Unit: $\text{J kg}^{-1}\ \text{K}^{-1}$.

Examples:

\begin{itemize}
\item water: $c \approx 4200\ \text{J kg}^{-1}\ \text{K}^{-1}$ (high — why water is a good coolant and why oceans steady the climate).

\item aluminium: $c \approx 900\ \text{J kg}^{-1}\ \text{K}^{-1}$.

\item copper: $c \approx 385\ \text{J kg}^{-1}\ \text{K}^{-1}$.

\end{itemize}

To find an unknown $c$ by experiment: supply known energy $Q$ electrically ($Q = VIt$, from the \textbf{power} 功率), then measure the temperature rise $\Delta T$ of a known \textbf{mass} 质量 $m$. Then $c = Q/(m\Delta T)$. Reduce heat loss with \textbf{insulation} 隔热 and use a rise of about 10 K (big enough to measure well, small enough to limit losses).

When two bodies reach thermal equilibrium with no heat lost to the surroundings, the energy gained by the colder one equals the energy lost by the hotter one:

\[
m_{1} c_{1} (T_{\text{eq}} - T_{1}) = m_{2} c_{2} (T_{2} - T_{\text{eq}}).
\]

\section*{Specific latent heat}

When a substance changes state (solid ↔ liquid, or liquid ↔ gas) at constant temperature, energy must be supplied (or removed) with \textbf{no} temperature change. This energy is the \textbf{latent heat} 潜热.

The \textbf{specific latent heat} 比潜热 $L$ is the energy to change the state of \textbf{unit mass} at constant temperature:

\[
L = \frac{Q}{m} \qquad\Longleftrightarrow\qquad Q = m L.
\]

Unit: $\text{J kg}^{-1}$.

Two kinds:

\begin{itemize}
\item \textbf{specific latent heat of fusion} 熔化 $L_{\text{f}}$ — for melting or freezing (solid ↔ liquid).

\item \textbf{specific latent heat of vaporisation} 汽化 $L_{\text{v}}$ — for boiling or condensing (liquid ↔ gas).

\end{itemize}

For water at atmospheric pressure: $L_{\text{f}} \approx 3.34 \times 10^{5}\ \text{J kg}^{-1}$ (at $0\ ^{\circ}\text{C}$); $L_{\text{v}} \approx 2.26 \times 10^{6}\ \text{J kg}^{-1}$ (at $100\ ^{\circ}\text{C}$). So $L_{\text{v}}$ is about 7 times $L_{\text{f}}$.

\subsection*{Why $L_{\text{v}} > L_{\text{f}}$}

Two reasons, both from the particle picture of matter:

\begin{enumerate}
\item \textbf{Bonds}: in melting, only some of the \textbf{intermolecular} 分子间 bonds break; the particles stay close as a liquid. In boiling, \textbf{all} the bonds must break so the particles can separate. Breaking all of them needs more energy.

\item \textbf{Work against the atmosphere}: when a liquid turns to gas it expands hugely (vapour has about $10^{3}$ times the liquid's \textbf{volume} 体积), so it does work pushing back the surrounding \textbf{atmospheric pressure} 大气压强. That work comes from the energy supplied.

\end{enumerate}

\subsection*{Multi-step problems}

If a problem mixes temperature change and a \textbf{phase change} 相变 (e.g. ice at $-5\ ^{\circ}\text{C}$ warming to water at $30\ ^{\circ}\text{C}$):

\begin{enumerate}
\item heat the solid from $-5\ ^{\circ}\text{C}$ to $0\ ^{\circ}\text{C}$: $Q_{1} = m c_{\text{ice}} \times 5$.

\item melt at $0\ ^{\circ}\text{C}$: $Q_{2} = m L_{\text{f}}$.

\item heat the water from $0\ ^{\circ}\text{C}$ to $30\ ^{\circ}\text{C}$: $Q_{3} = m c_{\text{water}} \times 30$.

\end{enumerate}

Total: $Q_{1} + Q_{2} + Q_{3}$. A phase change is at constant temperature, so use $mL$ there, not $mc\Delta T$.

When a question gives heater power $P$ and asks for the time, use $Q = Pt$ (assuming no heat loss). Insulating (lagging) the container and using a small mass are common ways to improve the experiment.


\end{document}
